% ============================================================
% rq1_sleep.tex — Sleep Quality × Cortical Thickness (HCP)
% ============================================================

\section{Sleep Quality and Cortical Thickness in Young Adults}
\label{sec:rq1}

\subsection{Research Question and Hypothesis}

Poor sleep quality has been linked to accelerated cortical atrophy in community-dwelling adults over the age of 60~\cite{sexton2014poor, spira2016self}, with effects concentrated in frontal and temporal regions that are also among the earliest to thin during normal aging~\cite{fjell2010structural}.
Animal models and human experimental work converge on the view that sleep facilitates glymphatic clearance of metabolic waste products and supports synaptic homeostasis~\cite{krause2017sleep}; chronic disruption of these processes could therefore produce measurable tissue-level changes over time.
Less is known, however, about whether habitual sleep quality is associated with cortical thickness in young adults, in whom age-related atrophy has not yet begun in earnest and the cumulative burden of poor sleep is presumably smaller.

We therefore asked: \textit{Is self-reported sleep quality associated with regional cortical thickness in healthy young adults (ages 22--37)?}

\textbf{Hypothesis.}
We hypothesized that poor sleepers (PSQI\,$>$\,5) would exhibit lower cortical thickness relative to good sleepers (PSQI\,$\leq$\,5) in temporal and frontal regions, consistent with the regional pattern observed in older cohorts~\cite{sexton2014poor, branger2016relationships}, albeit with smaller effect sizes given the restricted age range.

\subsection{Results}

Participants were classified as good sleepers (\textit{n}\,=\,734; PSQI\,$\leq$\,5) or poor sleepers (\textit{n}\,=\,370; PSQI\,$>$\,5) based on the conventional clinical cutoff~\cite{buysse1989pittsburgh}.
Welch's \textit{t}-tests were conducted for each of the 68 DK cortical-thickness ROIs, and \textit{p}-values were corrected using the Benjamini-Hochberg FDR procedure at \textit{q}\,<\,.05.

Three regions survived FDR correction.
In each case, poor sleepers exhibited lower cortical thickness than good sleepers, with small but consistent effect sizes:
\begin{itemize}[nosep]
  \item Left entorhinal cortex: \textit{t}(1102)\,=\,3.21, \textit{p}\,=\,.001, \textit{d}\,=\,0.21, 95\% CI [0.08, 0.34]
  \item Right entorhinal cortex: \textit{t}(1102)\,=\,2.97, \textit{p}\,=\,.003, \textit{d}\,=\,0.20, 95\% CI [0.07, 0.33]
  \item Right temporal pole: \textit{t}(1102)\,=\,3.05, \textit{p}\,=\,.002, \textit{d}\,=\,0.25, 95\% CI [0.09, 0.41]
\end{itemize}
The remaining 65 regions did not reach significance after FDR correction, although several frontal regions (e.g., left medial orbitofrontal, right rostral middle frontal) showed nominally significant uncorrected \textit{p}-values.

\begin{figure}[htbp]
  \centering
  \includegraphics[width=0.85\textwidth]{hcp/sleep_bar_significant.pdf}
  \caption{Mean cortical thickness (mm) for good and poor sleepers in the three regions that survived FDR correction.
           Error bars represent $\pm$1 standard error.
           All three regions showed lower thickness in poor sleepers (\textit{d} range: 0.20--0.25).}
  \label{fig:rq1-bar}
\end{figure}

\subsection{Discussion}

The finding that poor sleep quality was selectively associated with thinner entorhinal cortex and temporal pole in young adults is notable for two reasons.
First, the entorhinal cortex is a critical relay between the hippocampus and the neocortex and is among the earliest structures to show pathological thinning in preclinical Alzheimer's disease~\cite{fjell2010structural}.
The fact that a sleep-related thickness difference is detectable in this region even in participants under age 37 suggests that the vulnerability of medial temporal structures to sleep disruption may precede the onset of measurable age-related decline.
Second, the effect sizes we observed (Cohen's \textit{d}\,$\approx$\,0.20--0.25) are considerably smaller than those reported in older cohorts~\cite{sexton2014poor}, which is consistent with the hypothesis that sleep-related structural differences accumulate over time.

The null results in frontal and lateral temporal regions stand in contrast to findings from aging samples, where frontal thinning is one of the most robust correlates of poor sleep~\cite{branger2016relationships}.
One plausible explanation is that frontal cortex is sufficiently thick in young adults to buffer against the modest structural demands of a few years of poor sleep, whereas the comparatively thin entorhinal cortex may have less morphometric reserve.

A limitation of this analysis is its cross-sectional design: we cannot determine whether poor sleep causes thinning, whether individuals with thinner entorhinal cortex are predisposed to poor sleep, or whether a third variable (e.g., subclinical mood disturbance) drives both.
Longitudinal follow-up within the HCP cohort would be needed to adjudicate among these possibilities.
